\documentclass[../main.tex]{subfiles}

\begin{document}
\section{Data}

\subsection{Description and Exploration}
\noindent
The dataset used for this assignment is the PathMNIST subset distributed on
Canvas, derived from the MedMNIST v2 benchmark of standardised biomedical
images~\cite{yang2023medmnistv2}. The original PathMNIST collection contains
over one hundred thousand non-overlapping 28$\times$28 colour patches taken
from haematoxylin-and-eosin-stained slides of human colorectal cancer
tissue~\cite{kather2019colorectal}. The Canvas subset preserves the nine
tissue classes of the original release and ships pre-split into 32{,}000
training and 8{,}000 test samples, all stored as \texttt{uint8} tensors of
shape $(N, 28, 28, 3)$.

\noindent
The nine classes follow the official MedMNIST v2 definition and cover both
benign (e.g.~\emph{Adipose}, \emph{Lymphocytes}, \emph{Normal Colon Mucosa})
and malignant (\emph{Colorectal Adenocarcinoma Epithelium}) tissue, plus
classes that capture acellular elements such as \emph{Mucus} and
\emph{Background}. Inspection of the per-class sample counts in the training
set gives a maximum-to-minimum ratio of approximately $1.7\!:\!1$, so the
data is mildly imbalanced but not severely so; we therefore decided not to
re-weight the loss or apply class-conditional re-sampling.

\noindent
A qualitative grid of three random images per class
(Figure~\ref{fig.samples}) confirms that several classes share a similar
pinkish stain and a granular texture
(e.g.~\emph{Cancer-associated Stroma} versus \emph{Smooth Muscle}, or
\emph{Debris} versus \emph{Mucus}). This visual overlap predicts that
classes with subtle textural differences will be the hardest cases for any
classifier that does not model spatial structure explicitly.

\begin{figure}[H]
    \centering
    \includegraphics[width=0.55\linewidth]{Photos/samples_per_class.png}
    \caption{Three random PathMNIST training samples per class. The shared
    eosin-pink palette and the textural overlap between
    \emph{Smooth Muscle}/\emph{Cancer-associated Stroma} and
    \emph{Debris}/\emph{Mucus} are visible.}
    \label{fig.samples}
\end{figure}

\noindent
Per-channel pixel statistics show that the green channel has a noticeably
lower mean ($\approx 136$) than the red ($\approx 189$) and blue ($\approx
180$) channels, with a standard deviation of $\approx 46$ versus $\approx
32$. This is consistent with the absorption profile of the eosin
counter-stain and is one reason why we keep all three colour channels: a
grayscale conversion would discard a strong, class-correlated signal.

\subsection{Pre-processing}
\noindent
Two pre-processing pipelines are produced because the three classifiers
consume the data in different shapes and have different sensitivities to
feature scale.

\begin{itemize}
    \item \textbf{Neural-network branch (\gls{MLP} and \gls{CNN}).} Pixel
    intensities are cast to \texttt{float32} and divided by 255, mapping
    each pixel into $[0, 1]$. This keeps activations bounded at the start
    of training and prevents any single channel from dominating the
    gradients by virtue of its scale~\cite{lecun2012efficient}.
    \item \textbf{Classical-ML branch (\gls{SVM}).} Each image is flattened
    from $(28, 28, 3)$ into a $2{,}352$-dimensional feature vector. The
    flattened vector is then standardised (zero mean, unit variance per
    feature) before being projected by \gls{PCA} as part of the model
    pipeline. Standardisation is necessary because the \gls{RBF} kernel is
    scale sensitive; \gls{PCA} controls the runtime of the kernel
    evaluation and removes the strongest pixel-level correlations.
\end{itemize}

\noindent
Augmentation (random flips, small rotations) was considered for the
neural-network branch but ultimately not applied. The PathMNIST images are
already drawn from random tile positions inside the source slides, so the
training set covers a wide range of orientations naturally; adding
augmentation on top of an already-imbalanced 32{,}000-sample dataset would
have lengthened the hyper-parameter search without an obvious upside. For
similar reasons we kept the original 28$\times$28 resolution rather than
upsampling.

\noindent
For hyper-parameter selection an additional 80/20 stratified split is
carved from the provided training set. The held-out 6{,}400-sample
validation partition is used to score every candidate hyper-parameter
combination; the provided 8{,}000-sample test set is reserved exclusively
for the final evaluation in Section~\ref{section.4}. Stratification preserves
the class proportions on both sides of the split, which we verified
empirically: each class differs by at most one percentage point between the
training and validation partitions.

\end{document}
